\documentclass[10pt,a4paper]{amsart}
\usepackage[margin=19mm]{geometry}
\usepackage{amsmath,amssymb,mathtools}
\usepackage[scr=rsfs,cal=euler]{mathalfa}
\usepackage{xfrac}
\usepackage[unicode,hidelinks,bookmarks=false]{hyperref}
\newcommand{\ie}{i.e.\ }
\newcommand{\cf}{cf.\ }
\newcommand{\resp}{respectively\ }
\newcommand{\Subsection}{\subsection}
\providecommand{\nobreakdash}{\nobreak\hbox{-}}
\setlength{\emergencystretch}{2em}
\multlinegap0pt
\usepackage{multirow}
\usepackage{mathrsfs}
\usepackage{appendix}
\usepackage{textcomp}
\usepackage{manyfoot}
\usepackage{booktabs}
\usepackage{algorithm}
\usepackage{algorithmicx}
\usepackage{algpseudocode}
\usepackage{listings}
\newtheorem{theorem}{Theorem}%  meant for continuous numbers
\newtheorem{proposition}[theorem]{Proposition}% 
\newtheorem{conjecture}{Conjecture}
\newtheorem{lemma}{Lemma}
\newtheorem{example}{Example}%
\newtheorem{remark}{Remark}%
\newtheorem{definition}{Definition}%
\newtheorem{corollary}{Corollary}
\begin{document}
\title[Minimal Binary Linear Codes of Dimension n+4 from Partial Sp]{Minimal Binary Linear Codes of Dimension n+4 from Partial Spreads and Their Dual Access Structures}
\author{Apurba Sarkar, Kalyan Hansda,; Makhan Maji}
\date{}
\maketitle
\noindent\emph{Source and licence.} Minimal Binary Linear Codes of Dimension n+4 from Partial Spreads and Their Dual Access Structures. Version arXiv:2608.04889v1 (CC BY 4.0). This material is distributed under the Creative Commons Attribution 4.0 International licence, http://creativecommons.org/licenses/by/4.0/, as recorded in the arXiv OAI licence metadata for this exact version and shown on the abstract page https://arxiv.org/abs/2608.04889v1. Full rights notice: licenses/arxiv-2608.04889-CC-BY-4.0.txt.\par\smallskip
\noindent\textit{Editorial scope.} Counted unit: the original \texttt{theorem} environment \texttt{Thm:SSS\_Construction} at source lines 759--769 together with its complete proof at lines 771--775; the delivered document keeps source lines 183--776 so that every referenced label and every citation resolves inside it. Native editable source is the paper's own concatenated TeX; the AMS wrapper is editorial formatting only. No publisher class, upstream script or unrelated artwork is redistributed.
\begin{equation}\label{code}
\mathcal{M}_{(\psi_a)} = \left\{ m_{(\psi_a)}(\mathbf{\beta}) = \left( \sum_{k=1}^{4} \alpha_k \psi_k(\mathbf{x}) + \mathbf{\beta} \cdot \mathbf{x} \right)_{\mathbf{x} \in \mathbb{F}_2^{n*}} \;\middle|\; \alpha_k \in \mathbb{F}_2,\; \mathbf{\beta} \in \mathbb{F}_2^n \right\}.
\end{equation}

\begin{theorem}\label{C_weight}
Let $\mathcal{M}_{(\psi_a)}$ be the binary code defined in \eqref{code}. If $\phi(\mathbf{x})\neq \mathbf{v}\cdot \mathbf{x}$ for all $\mathbf{v} \in \mathbb{F}_2^n$ and all $\phi \in \mathcal{F}$, then $\mathcal{M}_{(\psi_a)}$ is a $[2^n-1, n+4]$ binary linear code whose weight distribution is given by the multiset union:
\begin{equation}\label{eq:multiset_weight}
\Omega = \bigcup_{\phi \in \mathcal{F}} \left\{ \frac{2^{n} - \hat{\phi}(\mathbf{\beta})}{2} \;\middle|\; \mathbf{\beta} \in \mathbb{F}_{2}^{n} \right\} \;\cup\; \{2^{n-1} \mid \mathbf{\beta} \in \mathbb{F}_{2}^{n*}\} \;\cup\; \{0\}.
\end{equation}
\end{theorem}

\begin{proof}
The length of $\mathcal{M}_{(\psi_a)}$ is directly determined by the cardinality of $\mathbb{F}_2^{n*}$, which is $2^n - 1$. Let $\{\mathbf{u}_j \mid j = 1,2,\ldots,n\}$ be a basis of the vector space $\mathbb{F}_2^n$. To establish the dimension, we show that the set $\mathcal{B} = \{m_{(\psi_a)}(\mathbf{u}_j) \mid 1 \leq j \leq n\} \cup \{m_{\psi_a}(\mathbf{0}) \mid 1 \leq a \leq 4\}$ forms a basis for $\mathcal{M}_{(\psi_a)}$. Setting a linear combination to zero gives:
\begin{align*}
&\sum_{j=1}^{n} \delta_j m_{(\psi_a)}(\mathbf{u}_j) + \sum_{k=1}^{4} \delta_{n+k} m_{\psi_k}(\mathbf{0}) = \mathbf{0}\\
\Rightarrow & \sum_{j=1}^{n} \delta_j 
\left(
\sum_{k=1}^{4} \psi_k(\mathbf{x}) + \mathbf{u}_j \cdot \mathbf{x}
\right)_{\mathbf{x} \in \mathbb{F}_2^{n*}}
+ \sum_{k=1}^{4} \delta_{n+k} (\psi_k(\mathbf{x}))_{\mathbf{x} \in \mathbb{F}_2^{n*}} = \mathbf{0}\\
\Rightarrow &\sum_{j=1}^{n} \delta_j (\mathbf{u}_j \cdot \mathbf{x}) + \sum_{k=1}^{4} \gamma_k \psi_k(\mathbf{x}) = 0, 
\end{align*}
where $\gamma_k = \sum_{j=1}^{n} \delta_j + \delta_{n+k}$. If $\gamma_k \neq 0$ for any $k$, it implies that some linear combination $\phi \in \mathcal{F}$ equals an inner product $\mathbf{v} \cdot \mathbf{x}$, which contradicts our initial assumption. Thus, $\gamma_k = 0$ for all $k$, which directly forces $\delta_j = 0$ for all $1 \leq j \leq n+4$. Since no subset of size $n+5$ can be linearly independent within this construction, the dimension of the code is exactly $n+4$. 

To establish the weight distribution, we evaluate the Hamming weight of the codewords using the Walsh-Hadamard transform. For the case where $\alpha_k = 1$ for all $k = 1,2,3,4$, the codeword corresponds to the evaluation vector of the full functional sum $\sum_{k=1}^{4} \psi_k$. Its Walsh-Hadamard transform at a shift vector $\mathbf{\beta} \in \mathbb{F}_2^n$ is given by:
\begin{align*}
\widehat{\sum_{k=1}^{4} \psi_k}(\mathbf{\beta}) 
&= \sum_{\mathbf{x} \in \mathbb{F}_2^n} (-1)^{\sum_{k=1}^{4} \psi_k(\mathbf{x}) + \mathbf{\beta} \cdot \mathbf{x}} \\
&= \left(2^n - \mathrm{wt}(m_{(\psi_a)}(\mathbf{\beta}))\right) - \mathrm{wt}(m_{(\psi_a)}(\mathbf{\beta})) \\
&= 2^n - 2\,\mathrm{wt}(m_{(\psi_a)}(\mathbf{\beta})).
\end{align*}
Isolating the weight term yields the explicit coordinate weight expression:
\[
\mathrm{wt}(m_{(\psi_a)}(\mathbf{\beta})) = \frac{2^n - \widehat{\sum_{k=1}^{4} \psi_k}(\mathbf{\beta})}{2}.
\]
The weight enumerations for all remaining configurations of the coefficients $\alpha_k \in \mathbb{F}_2$ are derived in an exactly analogous manner.
\end{proof}

We now establish a necessary and sufficient condition for the linear code $\mathcal{M}_{(\psi_a)}$ to achieve minimality based on its Walsh-Hadamard transform coefficients.
\begin{theorem}\label{Thm_Minimal_Cond}
The binary linear code $\mathcal{M}_{(\psi_a)}$ defined in \eqref{code} is minimal if and only if the following two conditions hold simultaneously for all vectors $\mathbf{x},\mathbf{y} \in \mathbb{F}_2^n$:
\begin{enumerate}
    \item For any $\phi_1,\phi_2 \in \mathcal{F}$ and $\mathbf{x} \neq \mathbf{y}$:
    \[ \widehat{\phi_1}(\mathbf{x})+\widehat{\phi_2}(\mathbf{y})\neq 2^n \quad \text{and} \quad \widehat{\phi_1}(\mathbf{x})-\widehat{\phi_2}(\mathbf{y})\neq 2^n. \tag{2.3} \]
    \item For any distinct $\phi_1,\phi_2 \in \mathcal{F}$:
    \[ \widehat{\phi_1}(\mathbf{x}+\mathbf{y})+\widehat{\phi_2}(\mathbf{x})-\widehat{\phi_1+\phi_2}(\mathbf{y})\neq 2^n. \tag{2.4} \]
\end{enumerate}
\end{theorem}

\begin{proof}
Every codeword $\mathbf{m}_{(\psi_a)}(\mathbf{\beta}) \in \mathcal{M}_{(\psi_a)}$ can be uniquely decomposed as follows:
\[
m_{(\psi_a)}(\mathbf{\beta})=\sum_{k=1}^{4}\alpha_k \mathbf{\psi}_k + c(\mathbf{\beta}),
\]
where $\mathbf{\psi}_a = \bigl(\psi_a(\mathbf{x})\bigr)_{\mathbf{x} \in \mathbb{F}_2^{n*}}$ and $c(\mathbf{\beta})$ is an element of the subcode:
\[
\mathcal{C}=\left\{c(\mathbf{\beta})=\bigl(\mathbf{\beta}\cdot \mathbf{x}\bigr)_{\mathbf{x} \in \mathbb{F}_2^{n*}} \mid \mathbf{\beta}\in \mathbb{F}_2^n\right\}.
\]
Here, $\mathcal{C}$ is a simplex code with parameters $[2^n - 1, n, 2^{n-1}]$. Thus, for any non-zero linear shift vector, $\mathrm{wt}(c(\mathbf{\beta})) = 2^{n-1}$. 

We verify that no improper support containment occurs across all structural subcases:

\textbf{Case I:} Let both codewords share an identical linear shift vector component such that $\mathbf{m}_1, \mathbf{m}_2 \in \left\{\sum_{k=1}^{4} \alpha_k \mathbf{\psi}_k + c(\mathbf{\beta}) \mid \alpha_k \in \mathbb{F}_2 \right\}$. 

\textbf{(a)} Let $\mathbf{m}_1 = c(\mathbf{\beta})$ and $\mathbf{m}_2 = \sum_{k=1}^{4} \mathbf{\psi}_k + c(\mathbf{\beta})$. Utilizing the weight distributions established via Theorem~\ref{C_weight}, the coordinate ordering evaluates to:
\begin{align*}
\mathbf{m}_1 \preceq \mathbf{m}_2 &\Leftrightarrow 
\mathrm{wt}\!\left(\sum_{k=1}^{4} \mathbf{\psi}_k\right) 
= \mathrm{wt}\!\left(\sum_{k=1}^{4} \mathbf{\psi}_k + c(\mathbf{\beta})\right) - 2^{n-1}\\
& \Leftrightarrow \widehat{\sum_{k=1}^{4} \psi_k}(\mathbf{0}) - \widehat{\sum_{k=1}^{4} \psi_k}(\mathbf{\beta}) = 2^n.
\end{align*}
Testing the reverse ordering configuration yields:
\[
\mathbf{m}_2 \preceq \mathbf{m}_1 \Leftrightarrow 
\mathrm{wt}\!\left(\sum_{k=1}^{4} \mathbf{\psi}_k\right) 
= 2^{n-1} - \mathrm{wt}\!\left(\sum_{k=1}^{4} \mathbf{\psi}_k + c(\mathbf{\beta})\right)
\Leftrightarrow \widehat{\sum_{k=1}^{4} \psi_k}(\mathbf{0}) + \widehat{\sum_{k=1}^{4} \psi_k}(\mathbf{\beta}) = 2^n.
\]

\textbf{(b)} Let $\mathbf{m}_1 = \mathbf{\psi}_a + \mathbf{\psi}_b + c(\mathbf{\beta})$ and $\mathbf{m}_2 = \mathbf{\psi}_c + \mathbf{\psi}_d + c(\mathbf{\beta})$ for $1 \le a,b,c,d \le 4$ with $a \neq b$ and $c \neq d$. The coordinate containments evaluate to:
\begin{align*}
&\mathbf{m}_1 \preceq \mathbf{m}_2 \Leftrightarrow 
\mathrm{wt}(\mathbf{\psi}_a + \mathbf{\psi}_b + \mathbf{\psi}_c + \mathbf{\psi}_d) \\
&= \mathrm{wt}(\mathbf{\psi}_c + \mathbf{\psi}_d + c(\mathbf{\beta})) - \mathrm{wt}(\mathbf{\psi}_a + \mathbf{\psi}_b + c(\mathbf{\beta}))\\
&\Leftrightarrow \widehat{\psi_a + \psi_b + \psi_c + \psi_d}(\mathbf{0}) + \widehat{\psi_a + \psi_b}(\mathbf{\beta}) - \widehat{\psi_c + \psi_d}(\mathbf{\beta}) = 2^n.
\end{align*}
Symmetrically, checking the opposite directional relation yields:
\[
\mathbf{m}_2 \preceq \mathbf{m}_1 \Leftrightarrow 
\mathrm{wt}(\mathbf{\psi}_a + \mathbf{\psi}_b + \mathbf{\psi}_c + \mathbf{\psi}_d) 
= \mathrm{wt}(\mathbf{\psi}_a + \mathbf{\psi}_b + c(\mathbf{\beta})) - \mathrm{wt}(\mathbf{\psi}_c + \mathbf{\psi}_d + c(\mathbf{\beta}))
\]
\[
\Leftrightarrow \widehat{\psi_a + \psi_b + \psi_c + \psi_d}(\mathbf{0}) + \widehat{\psi_c + \psi_d}(\mathbf{\beta}) - \widehat{\psi_a + \psi_b}(\mathbf{\beta}) = 2^n.
\]

\textbf{Case II:} Let the codewords be defined over distinct linear shift vectors such that $\mathbf{\beta}_1 \neq \mathbf{\beta}_2$:
\[
\mathbf{m}_1 \in \left\{\sum_{k=1}^{4} \alpha_k \mathbf{\psi}_k + c(\mathbf{\beta}_1) \;\middle|\; \alpha_k \in \mathbb{F}_2 \right\} \quad \text{and} \quad \mathbf{m}_2 \in \left\{\sum_{k=1}^{4} \alpha_k \mathbf{\psi}_k + c(\mathbf{\beta}_2) \;\middle|\; \alpha_k \in \mathbb{F}_2 \right\}.
\]

\textbf{(a)} Let $\mathbf{m}_1 = c(\mathbf{\beta}_1)$ and $\mathbf{m}_2 = \sum_{k=1}^{4} \mathbf{\psi}_k + c(\mathbf{\beta}_2)$. Evaluating the translations yields:
\[
\mathbf{m}_1 \preceq \mathbf{m}_2 \Leftrightarrow 
\mathrm{wt}\!\left(\sum_{k=1}^{4} \mathbf{\psi}_k + c(\mathbf{\beta}_1 + \mathbf{\beta}_2)\right)
= \mathrm{wt}\!\left(\sum_{k=1}^{4} \mathbf{\psi}_k + c(\mathbf{\beta}_2)\right) - 2^{n-1}
\]
\[
\Leftrightarrow \widehat{\sum_{k=1}^{4} \psi_k}(\mathbf{\beta}_1 + \mathbf{\beta}_2) - \widehat{\sum_{k=1}^{4} \psi_k}(\mathbf{\beta}_2) = 2^n.
\]
Reversing the containment evaluation gives:
\[
\mathbf{m}_2 \preceq \mathbf{m}_1 \Leftrightarrow 
\mathrm{wt}\!\left(\sum_{k=1}^{4} \mathbf{\psi}_k + c(\mathbf{\beta}_1 + \mathbf{\beta}_2)\right)
= 2^{n-1} - \mathrm{wt}\!\left(\sum_{k=1}^{4} \mathbf{\psi}_k + c(\mathbf{\beta}_2)\right)
\]
\[
\Leftrightarrow \widehat{\sum_{k=1}^{4} \psi_k}(\mathbf{\beta}_1 + \mathbf{\beta}_2) + \widehat{\sum_{k=1}^{4} \psi_k}(\mathbf{\beta}_2) = 2^n.
\]

\textbf{(b)} Let $\mathbf{m}_1 = \boldsymbol{\phi} + c(\mathbf{\beta}_1)$ and $\mathbf{m}_2 = \boldsymbol{\phi} + c(\mathbf{\beta}_2)$ with $\boldsymbol{\phi} = \mathbf{\psi}_a + \mathbf{\psi}_b + \mathbf{\psi}_c$ for $1\leq a,b,c \leq 4$ and $a \neq b \neq c$. Evaluating the vector shifts maps to:
\[
\mathbf{m}_1 \preceq \mathbf{m}_2 \Leftrightarrow 
\mathrm{wt}(c(\mathbf{\beta}_1 + \mathbf{\beta}_2)) 
= \mathrm{wt}(\mathbf{\psi}_a + \mathbf{\psi}_b + \mathbf{\psi}_c + c(\mathbf{\beta}_2)) - \mathrm{wt}(\mathbf{\psi}_a + \mathbf{\psi}_b + \mathbf{\psi}_c + c(\mathbf{\beta}_1))
\]
\[
\Leftrightarrow \widehat{\psi_a + \psi_b + \psi_c}(\mathbf{\beta}_1) - \widehat{\psi_a + \psi_b + \psi_c}(\mathbf{\beta}_2) = 2^n.
\]
Evaluating the opposite structural direction yields:
\[
\mathbf{m}_2 \preceq \mathbf{m}_1 \Leftrightarrow 
\mathrm{wt}(c(\mathbf{\beta}_1 + \mathbf{\beta}_2)) 
= \mathrm{wt}(\mathbf{\psi}_a + \mathbf{\psi}_b + \mathbf{\psi}_c+ c(\mathbf{\beta}_1)) - \mathrm{wt}(\mathbf{\psi}_a + \mathbf{\psi}_b + \mathbf{\psi}_c + c(\mathbf{\beta}_2))
\]
\[
\Leftrightarrow \widehat{\psi_a + \psi_b + \psi_c}(\mathbf{\beta}_2) - \widehat{\psi_a + \psi_b + \psi_c}(\mathbf{\beta}_1) = 2^n.
\]

\textbf{(c)} Let $\mathbf{m}_1 = \boldsymbol{\phi}_1 + c(\mathbf{\beta}_1)$ and $\mathbf{m}_2 = \boldsymbol{\phi}_2 + c(\mathbf{\beta}_2)$ with $\boldsymbol{\phi}_1 = \mathbf{\psi}_a + \mathbf{\psi}_b$ and $\boldsymbol{\phi}_2 = \mathbf{\psi}_c + \mathbf{\psi}_d$ for distinct $\boldsymbol{\phi}_1 \neq \boldsymbol{\phi}_2$. Resolving the underlying code constraints gives:
\[
\mathbf{m}_1 \preceq \mathbf{m}_2 \Leftrightarrow 
\mathrm{wt}\!\left(\sum_{k=1}^{4} \mathbf{\psi}_k + c(\mathbf{\beta}_1 + \mathbf{\beta}_2)\right) 
= \mathrm{wt}(\mathbf{\psi}_c + \mathbf{\psi}_d + c(\mathbf{\beta}_2)) - \mathrm{wt}(\mathbf{\psi}_a + \mathbf{\psi}_b + c(\mathbf{\beta}_1))
\]
\[
\Leftrightarrow \widehat{\sum_{k=1}^{4} \psi_k}(\mathbf{\beta}_1 + \mathbf{\beta}_2) + \widehat{\psi_a + \psi_b}(\mathbf{\beta}_1) - \widehat{\psi_c + \psi_d}(\mathbf{\beta}_2) = 2^n,
\]
and symmetrically for the inverse containment boundary:
\[
\mathbf{m}_2 \preceq \mathbf{m}_1 \Leftrightarrow 
\mathrm{wt}\!\left(\sum_{k=1}^{4} \mathbf{\psi}_k + c(\mathbf{\beta}_1 + \mathbf{\beta}_2)\right) 
= \mathrm{wt}(\mathbf{\psi}_a + \mathbf{\psi}_b + c(\mathbf{\beta}_1)) - \mathrm{wt}(\mathbf{\psi}_c + \mathbf{\psi}_d + c(\mathbf{\beta}_2))
\]
\[
\Leftrightarrow\widehat{\sum_{k=1}^{4} \psi_k}(\mathbf{\beta}_1 + \mathbf{\beta}_2) +   \widehat{\psi_c + \psi_d}(\mathbf{\beta}_2) - \widehat{\psi_a + \psi_b}(\mathbf{\beta}_1) = 2^n.
\]
Because none of these improper algebraic equalities can occur under the simultaneous constraints of equations (2.3) and (2.4), every non-zero codeword is minimal.
\end{proof}




\section{A family of minimal binary linear codes with dimension $n+4$}

 Let $n$ be a positive even integer and let $m=\frac{n}{2}$. A partial spread of order $\kappa$ in $\mathbb{F}_2^n$ is a collection of $m$-dimensional subspaces $\{V_1, V_2, \ldots, V_\kappa\}$ such that $V_a \cap V_b = \{\mathbf{0}\}$ for all $a \neq b$. Note that, $\kappa\leq 2^{t}+1$.

For each $1 \leq \ell \leq \kappa$, we define an Boolean function $\eta_\ell$ on $\mathbb{F}_2^n$ as:
\begin{equation}\label{eq:g_def}
\eta_\ell(\mathbf{x})=
\begin{cases}
1, & \text{if } \mathbf{x} \in V_\ell \setminus \{\mathbf{0}\},\\
0, & \text{otherwise.}
\end{cases} \tag{3.1}
\end{equation}
Let $\mathcal{S}_1, \mathcal{S}_2, \mathcal{S}_3, \mathcal{S}_4$ be distinct, non-empty subsets of $\{1,2,\ldots,\kappa\}$, with respective cardinalities $|\mathcal{S}_a| = s_a$. We define our boolean functions $\psi_a$ ($1 \leq a \leq 4$) as:
\begin{equation}\label{eq:f_def}
\psi_a = \sum_{\ell \in \mathcal{S}_a} \eta_\ell. \tag{3.2}
\end{equation}
It follows from the disjoint nature of the partial spread elements that any linear combination $\phi \in \mathcal{F}$ can be represented cleanly via symmetric differences of the index sets:
\begin{equation}\label{eq:f_comb}
\psi_a+\psi_b=\sum_{\ell\in \mathcal{S}_a\triangle \mathcal{S}_b}\eta_\ell, \quad \psi_a+\psi_b+\psi_c=\sum_{\ell\in \mathcal{S}_a\triangle \mathcal{S}_b\triangle \mathcal{S}_c}\eta_\ell, \quad \text{and} \quad \sum_{k=1}^{4}\psi_k=\sum_{\ell\in \mathcal{S}_1\triangle \mathcal{S}_2\triangle \mathcal{S}_3\triangle \mathcal{S}_4}\eta_\ell. \tag{3.3}
\end{equation}
We denote the cardinalities of these symmetric differences by $\sigma_{ab} = |\mathcal{S}_a \triangle \mathcal{S}_b|$, $\sigma_{abc} = |\mathcal{S}_a \triangle \mathcal{S}_b \triangle \mathcal{S}_c|$, and $\sigma_{1234} = |\mathcal{S}_1 \triangle \mathcal{S}_2 \triangle \mathcal{S}_3 \triangle \mathcal{S}_4|$.

\begin{theorem}\label{Thm_Spread_Weight}
Let the  sets $\mathcal{S}_a$ be distinct and satisfy $\mathcal{S}_1 \triangle \mathcal{S}_2 \triangle \mathcal{S}_3 \triangle \mathcal{S}_4 \neq \emptyset$. Then, the code $\mathcal{M}_{(\psi_a)}$ generated by the functions in \eqref{eq:f_def} is a $[2^n-1, n+4]$ binary linear code whose explicit weight parameters and multiplicities are uniquely determined by Table 1.
\end{theorem}

\begin{proof}
Let $\mathbf{\beta} \in \mathbb{F}_2^n$. The subspaces $V_\ell$ for $1 \leq \ell \leq \kappa$ are pairwise disjoint and $\dim(V_\ell) = m$. For any fixed index $\ell$, the sub-sum character evaluates to:
\[
\sum_{\mathbf{x} \in V_\ell} (-1)^{\mathbf{\beta} \cdot \mathbf{x}} = 
\begin{cases} 
2^m, & \text{if } \mathbf{\beta} \in V_\ell^\perp, \\ 
0, & \text{if } \mathbf{\beta} \notin V_\ell^\perp. 
\end{cases}
\]
Since $\psi_a(\mathbf{x}) = \sum_{\ell \in \mathcal{S}_a} \eta_\ell(\mathbf{x})$, where $\eta_\ell(\mathbf{x}) = 1$ if $\mathbf{x} \in V_\ell \setminus \{\mathbf{0}\}$ and $0$ otherwise, we have:
\begin{align*}
\widehat{\psi_a}(\mathbf{\beta}) &= \sum_{\mathbf{x} \in \mathbb{F}_2^n} (-1)^{\psi_a(\mathbf{x}) + \mathbf{\beta} \cdot \mathbf{x}} \\
&= \sum_{\mathbf{x} \in \mathbb{F}_2^n \setminus \bigcup_{\ell \in \mathcal{S}_a} V_\ell} (-1)^{\mathbf{\beta} \cdot \mathbf{x}} + \sum_{\ell \in \mathcal{S}_a} \sum_{\mathbf{x} \in V_\ell \setminus \{\mathbf{0}\}} (-1)^{1 + \mathbf{\beta} \cdot \mathbf{x}} + \sum_{\ell\in\mathcal{S}_a}(-1)^{\beta\cdot\mathbf{0}}\\
&= \sum_{\mathbf{x} \in \mathbb{F}_2^n} (-1)^{\mathbf{\beta} \cdot \mathbf{x}} - \sum_{\ell \in \mathcal{S}_a} \sum_{\mathbf{x} \in V_\ell} (-1)^{\mathbf{\beta} \cdot \mathbf{x}} - \sum_{\ell \in \mathcal{S}_a} \left[\sum_{\mathbf{x} \in V_\ell} (-1)^{\mathbf{\beta} \cdot \mathbf{x}} -(-1)^{\mathbf{\beta} \cdot \mathbf{0}} \right]+\sum_{\ell\in\mathcal{S}_a}(-1)^{\beta\cdot\mathbf{0}} \\
&= \sum_{\mathbf{x} \in \mathbb{F}_2^n} (-1)^{\mathbf{\beta} \cdot \mathbf{x}} - 2 \sum_{\ell \in \mathcal{S}_a} \sum_{\mathbf{x} \in V_\ell} (-1)^{\mathbf{\beta} \cdot \mathbf{x}} + 2\sum_{\ell \in \mathcal{S}_a} 1\\
&= \sum_{\mathbf{x} \in \mathbb{F}_2^n} (-1)^{\mathbf{\beta} \cdot \mathbf{x}} - 2 \sum_{\ell \in \mathcal{S}_a} \sum_{\mathbf{x} \in V_\ell} (-1)^{\mathbf{\beta} \cdot \mathbf{x}} + 2s_a.
\end{align*}

We establish the values of $\widehat{\psi_a}(\mathbf{\beta})$ by analyzing the orthogonal spaces across three cases:

\textbf{Case 1:} $\mathbf{\beta} = \mathbf{0}$.
\begin{align*}
\widehat{\psi_a}(\mathbf{0}) &= 2^n - 2 \sum_{\ell \in \mathcal{S}_a} 2^m + 2s_a \\
&= 2^n - 2s_a 2^m + 2s_a \\
&= 2^n - 2s_a(2^m - 1).
\end{align*}

\textbf{Case 2:} $\mathbf{\beta} \neq \mathbf{0}$ and $\mathbf{\beta} \notin V_\ell^\perp$ for all $\ell \in \mathcal{S}_a$.
Since $\mathbf{\beta} \neq \mathbf{0}$, t$\sum_{\mathbf{x} \in \mathbb{F}_2^n} (-1)^{\mathbf{\beta} \cdot \mathbf{x}} = 0$, and the subspace character sums vanish for all $\ell \in \mathcal{S}_a$. Thus:
\begin{align*}
\widehat{\psi_a}(\mathbf{\beta}) &= 0 - 2 \sum_{\ell \in \mathcal{S}_a} 0 + 2\sum_{\ell \in \mathcal{S}_a} 1 \\
&= 2s_a.
\end{align*}

\textbf{Case 3:} $\mathbf{\beta} \neq \mathbf{0}$ and $\mathbf{\beta} \in V_{\ell_0}^\perp$ for a unique index $\ell_0 \in \mathcal{S}_a$. 
Since $V_\ell \cap V_k = \{\mathbf{0}\}$ for $\ell \neq k$, it follows that $\dim(V_\ell + V_k) = 2m = n$, which guarantees $V_\ell^\perp \cap V_k^\perp = \{\mathbf{0}\}$. Thus, $\mathbf{\beta}$ can belong to at most one $V_{\ell_0}^\perp \setminus \{\mathbf{0}\}$.
\begin{align*}
\widehat{\psi_a}(\mathbf{\beta}) &= 0 - 2 \left( \sum_{\ell \in \mathcal{S}_a \setminus \{\ell_0\}} 0 + 2^m \right) + 2s_a \\
&= -2^{m+1} + 2s_a.
\end{align*}

Combining the cases yields the explicit system:
\[
\widehat{\psi_a}(\mathbf{\beta})=
\begin{cases}
2^n-2s_a(2^m-1), & \text{if } \mathbf{\beta}=\mathbf{0},\\
2s_a, & \text{if } \mathbf{\beta}\notin V_\ell^\perp,\ \forall \; \ell\in \mathcal{S}_a,\\
-2^{m+1}+2s_a, & \text{if } \mathbf{\beta}\in V_\ell^\perp\setminus\{\mathbf{0}\} \text{ for some } \ell\in \mathcal{S}_a.
\end{cases}
\]
Substituting $\widehat{\psi_a}(\mathbf{\beta})$ directly into the weight equation $\mathrm{wt}(m_{(\psi_a)}(\mathbf{\beta})) = \frac{2^n - \widehat{\psi_a}(\mathbf{\beta})}{2}$ maps out the unique multiplicities in Table 1.
\end{proof}

\begin{table}[ht]
\centering
\caption{The values of Walsh transforms of $\phi \in \mathcal{F}$}
\label{Table_1}
\label{tab:walsh}
\renewcommand{\arraystretch}{1.0}
\begin{tabular}{|l|l|l|l|}
\hline
Walsh transform &
$\mathbf{\beta=0}$ &
$\mathbf{\beta} \notin V_\ell^\perp,\ \forall\, \ell \in \mathcal{S}$ &
$\mathbf{\beta}\in V_\ell^\perp\setminus\{\mathbf{0}\}\; \text{for some }\ell \in \mathcal{S}$ \\
\hline
$\widehat{\psi_a}(\mathbf{\beta})$ &
$2^n-2s_a(2^m-1)$ &
$2s_a$ &
$-2^{m+1}+2s_a$
\\
\hline
$\widehat{\psi_a+\psi_b}(\mathbf{\beta})$ &
$2^n-2\sigma_{ab}(2^m-1)$ &
$2\sigma_{ab}$ &
$-2^{m+1}+2\sigma_{ab}$
\\
\hline
$\widehat{\psi_a+\psi_b+\psi_c}(\mathbf{\beta})$ &
$2^n-2\sigma_{abc}(2^m-1)$ &
$2\sigma_{abc}$ &
$-2^{m+1}+2\sigma_{abc}$
\\
\hline
$\widehat{\displaystyle\sum_{k=1}^{4}\psi_k}(\mathbf{\beta})$ &
$2^n-2\sigma_{1234}(2^m-1)$ &
$2\sigma_{1234}$ &
$-2^{m+1}+2\sigma_{1234}$
\\
\hline
\end{tabular}
\end{table}

To guarantee that the linear code $\mathcal{M}_{(\psi_a)}$ is structurally minimal, we establish the following geometric and combinatorial constraints on the underlying index sets:

\begin{itemize}
    \item[\textbf{C1:}] $\mathcal{S}_c \not\subset \mathcal{S}_a \triangle \mathcal{S}_b$, $\mathcal{S}_a \triangle \mathcal{S}_b \not\subset \mathcal{S}_c$, $\mathcal{S}_d \not\subset \mathcal{S}_a \triangle \mathcal{S}_b \triangle \mathcal{S}_c$, and $\mathcal{S}_a \triangle \mathcal{S}_b \triangle \mathcal{S}_c \not\subset \mathcal{S}_d$ for  $a,b,c,d \in \{1,2,3,4\}$.
    \item[\textbf{C2:}] $\bigcap_{a=1}^3 \mathcal{S}_a \neq \emptyset$, $\bigcap_{a=1}^{4} \mathcal{S}_a \neq \emptyset$, and the condition $|\mathcal{S}_a \cap \mathcal{S}_b \cap \mathcal{S}_c| \neq |\mathcal{S}_i \cap \mathcal{S}_j|$ holds for at least two pairs of distinct indices and the condition $|\mathcal{S}_a \cap \mathcal{S}_b \cap \mathcal{S}_c \cap \mathcal{S}_d| \neq |\mathcal{S}_i \cap \mathcal{S}_j \cap \mathcal{S}_k|$ holds for at least three pairs of distinct indices.
    \item[\textbf{C3:}] $\sigma_{ab} \geq 2$, $\sigma_{abc} \geq 2$, $\sigma_{ab} + \sigma_{ad} \neq \sigma_{bd}$, $\sigma_{ab} + \sigma_{acd} \neq \sigma_{bcd}$ and $\sigma_{abc} + \sigma_{abd} \neq \sigma_{cd}$.
\end{itemize}

These constraints generalize the structural criteria developed for lower-dimensional configurations in \cite{LiuLiao2022}.

\begin{table}[h]
\centering
\caption{Weight distribution of $\mathcal{M}_{(\psi_a)}$}
\label{Table_2}
\begin{tabular}{|c|c|}
\hline
Weight ($\omega$) & Multiplicity ($A_\omega$) \\
\hline
$0$ & $1$ \\
\hline
$u(2^m-1)$ & $1, \quad u \in \{s_a, \sigma_{ab}, \sigma_{abc}, \sigma_{1234}\}$ \\
\hline
$2^{n-1}$ & $2^n-1$ \\
\hline
$2^{n-1}-u$ & $(2^m+1-u)(2^m-1), \quad u \in \{s_a, \sigma_{ab}, \sigma_{abc}, \sigma_{1234}\}$ \\
\hline
$2^{n-1}+2^m-u$ & $u(2^m-1), \quad u \in \{s_a, \sigma_{ab}, \sigma_{abc}, \sigma_{1234}\}$ \\
\hline
\end{tabular}
\end{table}

We present two fundamental properties derived from these set conditions, which are required to establish the minimality proofs of the code space.
\begin{lemma}\label{lem:set_theory}
For $a \in \{1,2,3,4\}$, let $\mathcal{S}_a$ be non-empty index sets. If constraint C1 holds, then:
\begin{enumerate}
    \item $\mathcal{S}_a \not\subseteq \mathcal{S}_b$ and $\mathcal{S}_a \cap \mathcal{S}_b \neq \emptyset$ for all distinct $a, b \in \{1,2,3,4\}$.
    \item $(\mathcal{S}_a \cap \mathcal{S}_b) \triangle (\mathcal{S}_a \cap \mathcal{S}_c) \subsetneq \mathcal{S}_b \triangle \mathcal{S}_c$ for distinct $a,b,c$.
\end{enumerate}
\end{lemma}

\begin{proof}
\begin{enumerate}
    \item Let $\mathcal{S}_a \subseteq \mathcal{S}_b$ then  $\mathcal{S}_a \triangle \mathcal{S}_b = \mathcal{S}_b \setminus \mathcal{S}_a$. so that 
    $\mathcal{S}_b \setminus \mathcal{S}_a \subseteq \mathcal{S}_b$ implies that $\mathcal{S}_a \triangle \mathcal{S}_b \subseteq \mathcal{S}_b$. This contradicts  C1. 
    Let $\mathcal{S}_a \cap \mathcal{S}_b = \emptyset$.
    
    Now if  $\mathcal{S}_a \triangle \mathcal{S}_b = \mathcal{S}_a \cup \mathcal{S}_b$  we have  $\mathcal{S}_a \subseteq \mathcal{S}_a \triangle \mathcal{S}_b$. Again, contradicts  C1.
    Thus, $\mathcal{S}_a \not\subseteq \mathcal{S}_b$ and $\mathcal{S}_a \cap \mathcal{S}_b \neq \emptyset$.

    \item Let $x \in (\mathcal{S}_a \cap \mathcal{S}_b) \triangle (\mathcal{S}_a \cap \mathcal{S}_c)$. Then 
    \begin{align*}
    &x \in (\mathcal{S}_a \cap \mathcal{S}_b) \setminus (\mathcal{S}_a \cap \mathcal{S}_c)\\ 
    &\implies x \in \mathcal{S}_a, \, x \in \mathcal{S}_b, \, x \notin \mathcal{S}_c\\
    & \implies x \in \mathcal{S}_b \setminus \mathcal{S}_c \subseteq \mathcal{S}_b \triangle \mathcal{S}_c\\
    & \implies (\mathcal{S}_a \cap \mathcal{S}_b) \triangle (\mathcal{S}_a \cap \mathcal{S}_c) \subseteq \mathcal{S}_b \triangle \mathcal{S}_c.
    \end{align*}
    Assume $(\mathcal{S}_a \cap \mathcal{S}_b) \triangle (\mathcal{S}_a \cap \mathcal{S}_c) = \mathcal{S}_b \triangle \mathcal{S}_c$.
    Then $(\mathcal{S}_a \cap \mathcal{S}_b) \triangle (\mathcal{S}_a \cap \mathcal{S}_c) \subseteq \mathcal{S}_a$ implies that  $\mathcal{S}_b \triangle \mathcal{S}_c \subseteq \mathcal{S}_a$,  this contradicts  C1. Thus, $(\mathcal{S}_a \cap \mathcal{S}_b) \triangle (\mathcal{S}_a \cap \mathcal{S}_c) \subsetneq \mathcal{S}_b \triangle \mathcal{S}_c$.
\end{enumerate}
\end{proof}

\begin{lemma}\label{lem:bound_theory}
Let $\mathcal{S}_a$ be non-empty index sets satisfying $2 \leq s_a \leq 2^{m-1}$ for $a \in \{1,2,3,4\}$. If conditions C1 and C2 hold simultaneously, the symmetric difference targets are upper-bounded by:
\[
\sigma_{ab} \leq 2^m - 2, \quad \quad \sigma_{abc} \leq 2^m - 2 \quad \text{and} \quad \sigma_{abcd} \leq 2^m - 2.
\]
\end{lemma}

\begin{proof}
By  Lemma \ref{lem:set_theory}(1) we have that  $\sigma_{ab} = s_a + s_b - 2s_{ab}, \quad s_{ab} \geq 1$. 
Given $s_a, s_b \leq 2^{m-1}$ it yields that $\sigma_{ab} \leq 2^{m-1} + 2^{m-1} - 2(1) = 2^m - 2$. 

Now $|\mathcal{S}_1 \cup \mathcal{S}_2 \cup \mathcal{S}_3| \leq 2^m + 1$. 
Expanding it gives  $\sigma_{abc} = |\mathcal{S}_1 \cup \mathcal{S}_2 \cup \mathcal{S}_3| - \left( s_{12} + s_{13} + s_{23} - 2s_{123} \right)$.
By Condition C2, $s_{12} + s_{13} + s_{23} - 2s_{123} \geq 3$. Therefore $\sigma_{abc} \leq (2^m + 1) - 3 = 2^m - 2$. 

Now, $|\mathcal{S}_1 \cup \mathcal{S}_2 \cup \mathcal{S}_3 \cup \mathcal{S}_4| \leq 2^m + 1$. Expanding it gives:
\[
\sigma_{abcd} = |\mathcal{S}_1 \cup \mathcal{S}_2 \cup \mathcal{S}_3 \cup \mathcal{S}_4| - \left( \sum_{1 \le a < b \le 4} s_{ab} - 3 \sum_{1 \le a < b < c \le 4} s_{abc} + 7 s_{abcd} \right).
\]
By Condition C2, $\sum_{1 \le a < b \le 4} s_{ab} - 3 \sum_{1 \le a < b < c \le 4} s_{abc} + 7 s_{abcd} \ge 3$. Therefore, $\sigma_{abcd} \le (2^m + 1) - 3 = 2^m - 2$.
\end{proof}



By utilizing the bounded discrete transform values from Lemma \ref{lem:bound_theory}, we can consolidate the multi-case propositions into two overarching theorems that satisfy the minimality criteria of Theorem \ref{Thm_Minimal_Cond}.




\begin{theorem}\label{Thm_Verify_Cond1}
Let $n \geq 8$ and $2 \leq s_a \leq 2^{m-1}$ for all $1 \leq a \leq 4$. For any distinct vectors $\mathbf{x}, \mathbf{y} \in \mathbb{F}_2^n$, the Walsh-Hadamard combinations of the spread functions strictly satisfy:
\[ \widehat{\phi}_1(\mathbf{x}) + \widehat{\phi}_2(\mathbf{y}) \neq 2^n \quad \text{and} \quad \widehat{\phi}_1(\mathbf{x}) - \widehat{\phi}_2(\mathbf{y}) \neq 2^n, \]
for all choices of $\phi_1, \phi_2 \in \mathcal{F}$.
\end{theorem}

\begin{proof}
First, we have $\max_{\mathbf{\beta} \neq \mathbf{0}} |\widehat{\phi}(\mathbf{\beta})| \leq 2^{m+1} - 4$. Let $\sigma_1 = |\mathcal{S}_{\phi_1}|$ and $\sigma_2 = |\mathcal{S}_{\phi_2}|$. By Lemma \ref{lem:bound_theory}:
\begin{equation}\label{eq:thm_bounds}
2 \leq \sigma_1, \sigma_2 \leq 2^m - 2.
\end{equation}
We have the following cases:

\textbf{Case 1: $\mathbf{x} \neq \mathbf{0}$ and $\mathbf{y} \neq \mathbf{0}$.} \\
By triangle inequality, $|\widehat{\phi}_1(\mathbf{x}) \pm \widehat{\phi}_2(\mathbf{y})| \leq |\widehat{\phi}_1(\mathbf{x})| + |\widehat{\phi}_2(\mathbf{y})| \leq 2(2^{m+1}-4) = 2^{m+2}-8$, which  implies $\widehat{\phi}_1(\mathbf{x}) \pm \widehat{\phi}_2(\mathbf{y}) \neq 2^n$.


\textbf{Case 2: $\mathbf{x} = \mathbf{0}$ and $\mathbf{y} \neq \mathbf{0}$.} \\
Substituting $\mathbf{x} = \mathbf{0}$ into $\psi_a(\mathbf{x})=\sum_{\mathbf{x} \in \mathbb{F}_2^n} (-1)^{\mathbf{\beta} \cdot \mathbf{x}} - 2 \sum_{\ell \in \mathcal{S}_a} \sum_{\mathbf{x} \in V_\ell} (-1)^{\mathbf{\beta} \cdot \mathbf{x}} + 2s_a$  yields
that 
$\widehat{\phi}_1(\mathbf{0}) \pm \widehat{\phi}_2(\mathbf{y}) = 2^n - 2\sigma_1(2^m - 1) \pm \widehat{\phi}_2(\mathbf{y})$.
Now assume $\widehat{\phi}_1(\mathbf{0}) + \widehat{\phi}_2(\mathbf{y}) = 2^n$. This implies
\begin{equation}\label{eq:case2_pos_eq}
\widehat{\phi}_2(\mathbf{y}) = 2\sigma_1(2^m - 1).
\end{equation}
Since $\mathbf{y} \in V_{k_0}^\perp \setminus \{\mathbf{0}\}$ for at most one $k_0 \in \mathcal{S}_{\phi_2}$, we have 
\begin{equation}
\max_{\mathbf{y} \neq \mathbf{0}} \widehat{\phi}_2(\mathbf{y}) \leq 2\sigma_2 + 2^{m+1}. 
\end{equation}
Now, this together  with \eqref{eq:case2_pos_eq} gives 
$2\sigma_2 + 2^{m+1} \geq 2\sigma_1(2^m - 1)$ if and only if $\sigma_2 + 2^m \geq \sigma_1 2^m - \sigma_1$.

Applying \eqref{eq:thm_bounds} we have $(2^m - 2) + 2^m \geq 2^{m+1} - 2$ if and only if $(\sigma_1, \sigma_2) = (2, 2^m - 2)$.  By conditions C1 and C3, $\mathcal{S}_{\phi_1} \not\subseteq \mathcal{S}_{\phi_2}$ and $\mathcal{S}_{\phi_2} \not\subseteq \mathcal{S}_{\phi_1}$ implies $\sigma_2 \leq 2^m - 3$ when $\sigma_1 = 2$, which results in a contradiction.

Now assume $\widehat{\phi}_1(\mathbf{0}) - \widehat{\phi}_2(\mathbf{y}) = 2^n$. This implies that 
$-\widehat{\phi}_2(\mathbf{y}) = 2\sigma_1(2^m - 1)$.
Since $\sigma_1 \geq 2$, we have $2\sigma_1(2^m - 1) \geq 4(2^m - 1) = 2^{m+2} - 4$. However:
\[
\max_{\mathbf{y} \neq \mathbf{0}} |-\widehat{\phi}_2(\mathbf{y})| \leq 2^{m+1} - 2\sigma_2 < 2^{m+2} - 4,
\]
which is impossible.

By symmetry, the case $\mathbf{x} \neq \mathbf{0}, \mathbf{y} = \mathbf{0}$ holds identically. Thus, $\widehat{\phi}_1(\mathbf{x}) \pm \widehat{\phi}_2(\mathbf{y}) \neq 2^n$ for all $\mathbf{x} \neq \mathbf{y}$.
\end{proof}




\begin{theorem}\label{Thm_Verify_Cond2}
Let $n \geq 8$ and $2 \leq s_a \leq 2^{m-1}$ for all $1 \leq a \leq 4$. For all functional pairings $\phi_1, \phi_2 \in \mathcal{F}$ with $\phi_1 \neq \phi_2$, the shifting differences satisfy:
\[ \widehat{\phi_1}(\mathbf{x}+\mathbf{y})+\widehat{\phi_2}(\mathbf{x})-\widehat{\phi_1+\phi_2}(\mathbf{y})\neq 2^n, \]
across all vector spatial configurations of $\mathbf{x}, \mathbf{y} \in \mathbb{F}_2^n$.
\end{theorem}

\begin{proof}
Let $\phi_1, \phi_2 \in \mathcal{F}$ with $\phi_1 \neq \phi_2$, and define $\psi = \phi_1 + \phi_2 \in \mathcal{F}$. Let $\sigma_1 = |\mathcal{S}_{\phi_1}|$, $\sigma_2 = |\mathcal{S}_{\phi_2}|$, and $\sigma_3 = |\mathcal{S}_{\psi}| = |\mathcal{S}_{\phi_1} \triangle \mathcal{S}_{\phi_2}|$. For any $\mathbf{x}, \mathbf{y} \in \mathbb{F}_2^n$, define $\Lambda(\mathbf{x}, \mathbf{y}) = \widehat{\phi}_1(\mathbf{x}+\mathbf{y}) + \widehat{\phi}_2(\mathbf{x}) - \widehat{\psi}(\mathbf{y})$.


We evaluate $\Lambda(\mathbf{x}, \mathbf{y})$ across three disjoint spatial vector configurations:

\textbf{Case 1: $\mathbf{x} = \mathbf{y}$.} \\
When $\mathbf{x} = \mathbf{y}$,  $\Lambda(\mathbf{x}, \mathbf{x}) = \widehat{\phi}_1(\mathbf{0}) + \widehat{\phi}_2(\mathbf{x}) - \widehat{\psi}(\mathbf{x})$. Substituting the evaluation at the origin yields:
\[
\Lambda(\mathbf{x}, \mathbf{x}) = 2^n - 2\sigma_1(2^m - 1) + \widehat{\phi}_2(\mathbf{x}) - \widehat{\psi}(\mathbf{x}).
\]
Assuming $\Lambda(\mathbf{x}, \mathbf{x}) = 2^n$ requires $\widehat{\phi}_2(\mathbf{x}) - \widehat{\psi}(\mathbf{x}) = 2\sigma_1(2^m - 1)$. If $\mathbf{x} = \mathbf{0}$, this gives:
\[
2^n - 2\sigma_2(2^m - 1) - \left(2^n - 2\sigma_3(2^m - 1)\right) = 2\sigma_1(2^m - 1) \iff \sigma_3 = \sigma_1 + \sigma_2,
\]
which contradicts $|\mathcal{S}_{\phi_1} \triangle \mathcal{S}_{\phi_2}| \leq \sigma_1 + \sigma_2 - 2|\mathcal{S}_{\phi_1} \cap \mathcal{S}_{\phi_2}| < \sigma_1 + \sigma_2$ by condition C2 ($|\mathcal{S}_{\phi_1} \cap \mathcal{S}_{\phi_2}| \neq 0$). If $\mathbf{x} \neq \mathbf{0}$, then $\max_{\mathbf{x} \neq \mathbf{0}} |\widehat{\phi}_2(\mathbf{x}) - \widehat{\psi}(\mathbf{x})| \leq 2^{m+2} - 8$, whereas $2\sigma_1(2^m - 1) \geq 2^{m+2} - 4$, making equality impossible.

\textbf{Case 2: Either $\mathbf{x} = \mathbf{0}$ or $\mathbf{y} = \mathbf{0}$ with $\mathbf{x} \neq \mathbf{y}$.} \\
If $\mathbf{x} = \mathbf{0}$ and $\mathbf{y} \neq \mathbf{0}$, then $\Lambda(\mathbf{0}, \mathbf{y}) = \widehat{\phi}_1(\mathbf{y}) + \widehat{\phi}_2(\mathbf{0}) - \widehat{\psi}(\mathbf{y})$. Assume $\Lambda(\mathbf{0}, \mathbf{y}) = 2^n$:
\[
\widehat{\phi}_1(\mathbf{y}) - \widehat{\psi}(\mathbf{y}) = 2\sigma_2(2^m - 1).
\]
By the spectral triangle inequality:
\[
\max_{\mathbf{y} \neq \mathbf{0}} |\widehat{\phi}_1(\mathbf{y}) - \widehat{\psi}(\mathbf{y})| \leq (2\sigma_1 + 2^{m+1}) + (2\sigma_3 + 2^{m+1}) = 2(\sigma_1 + \sigma_3) + 2^{m+2}.
\]
This forces $2(\sigma_1 + \sigma_3) + 2^{m+2} \geq 2\sigma_2(2^m - 1)$, which yields a direct contradiction under constraints C1–C3. The symmetric case $\mathbf{y} = \mathbf{0}, \mathbf{x} \neq \mathbf{0}$ holds identically.

\textbf{Case 3: $\mathbf{x} \neq \mathbf{0}$, $\mathbf{y} \neq \mathbf{0}$, and $\mathbf{x} \neq \mathbf{y}$.} \\
Since $\mathbf{x}, \mathbf{y}, \mathbf{x}+\mathbf{y}$ are all non-zero, each vector belongs to at most one orthogonal subspace component $V_\ell^\perp$. By the triangle inequality:
\[
|\Lambda(\mathbf{x}, \mathbf{y})| \leq |\widehat{\phi}_1(\mathbf{x}+\mathbf{y})| + |\widehat{\phi}_2(\mathbf{x})| + |\widehat{\psi}(\mathbf{y})| \leq 3(2^{m+1} - 4) = 3 \cdot 2^{m+1} - 12.
\]
For $n \geq 8 \implies m \geq 4$:
\[
3 \cdot 2^{m+1} - 12 < 2^{2m} = 2^n.
\]
Thus, $\Lambda(\mathbf{x}, \mathbf{y}) \neq 2^n$ across all vector spatial configurations.
\end{proof}
\begin{theorem}\label{Thm_Main_Code_Result}
Let $n \ge 8$ be an even integer with $m = n/2$, and let the index set cardinalities satisfy $2 \le s_a \le 2^{m-1}$ for all $1 \le a \le 4$. Define the minimum partial spread intersection parameter across all 15 non-zero functions $\phi \in \mathcal{F}$ as:
\[
\theta = \min_{1 \le a,b,c \le 4} \{ s_a, \, \sigma_{ab}, \, \sigma_{abc}, \, \sigma_{1234} \}.
\]
If $\theta \le 2^{m-1} - 1$, then $\mathcal{M}_{(\psi_a)}$ is a valid $[2^n - 1, \, n + 4, \, \theta(2^m - 1)]$ minimal binary linear code. Furthermore, if $\theta \le 2^{m-2}$, the code strictly violates the Ashikhmin-Barg condition, exhibiting a Hamming weight ratio of:
\[
\frac{w_{\min}}{w_{\max}} \le \frac{1}{2}.
\]
\end{theorem}
\begin{proof}
By Theorems \ref{Thm_Verify_Cond1} and \ref{Thm_Verify_Cond2}, all spatial vector conditions required by Theorem \ref{Thm_Minimal_Cond} are simultaneously satisfied, establishing that every non-zero codeword in $\mathcal{M}_{(\psi_a)}$ is minimal.

From the weight spectrum derived in Table~\ref{Table_2}, the non-zero Hamming weights of $\mathcal{M}_{(\psi_a)}$ belong to the set:
\[
\Omega \setminus \{0\} = \{ u(2^m - 1), \; 2^{n-1}, \; 2^{n-1} - u, \; 2^{n-1} + 2^m - u \},
\]
where $u \in \{s_a, \sigma_{ab}, \sigma_{abc}, \sigma_{1234}\}$. The absolute minimum weight of $\mathcal{M}_{(\psi_a)}$ is attained at $u = \theta$, giving $w_{\min} = \theta(2^m - 1)$.

Conversely, the maximum weight is attained at $w_{\max} = 2^{n-1} + 2^m - \theta$. Evaluating the ratio $w_{\min} / w_{\max}$ under the structural threshold $\theta \le 2^{m-2}$ gives:
\[
\frac{w_{\min}}{w_{\max}} = \frac{\theta(2^m - 1)}{2^{n-1} + 2^m - \theta}.
\]
Since $n \ge 8$ implies $m \ge 4$, substituting $\theta \le 2^{m-2}$ yields:
\[
\frac{w_{\min}}{w_{\max}} \le \frac{2^{m-2}(2^m - 1)}{2^{2m-1} + 2^m - 2^{m-2}} = \frac{2^{2m-2} - 2^{m-2}}{2^{2m-1} + 3 \cdot 2^{m-2}} < \frac{1}{2}.
\]
Thus, the code strictly violates the Ashikhmin--Barg ceiling ($w_{\min} / w_{\max} \le 1/2$) while remaining strictly minimal.
\end{proof}

\begin{example}\label{ex:AB_strictly_greater}
Let $n = 8$ ($m = 4$), yielding code length $N = 2^8 - 1 = 255$ and dimension $K = n + 4 = 12$. We select $17$ pairwise disjoint $4$-dimensional subspaces $\{V_1, V_2, \ldots, V_{17}\}$ in $\mathbb{F}_2^8$. To ensure conditions C1--C3 and the size bounds $2 \le s_a \le 2^{m-1} = 8$ are strictly satisfied, we define the four basis index sets as:
\[
\mathcal{S}_1 = \{1, 2, 3, 4, 10, 11, 13\}, \quad \mathcal{S}_2 = \{3, 4, 5, 6, 10, 11, 13\},
\]
\[
\mathcal{S}_3 = \{1, 3, 7, 8, 10, 12, 13\}, \quad \text{and} \quad \mathcal{S}_4 = \{2, 3, 5, 9, 11, 12, 13\},
\]
giving individual set cardinalities $s_1 = s_2 = s_3 = s_4 = 7 \le 2^{4-1}$.
Evaluating all $15$ non-zero function combinations in $\mathcal{F}$ yields the parameter set:
\[
u \in \{\sigma_{ab}, \sigma_{abc}, \sigma_{1234}\} \subset \{5, 6, 7, 8\},
\]
thereby determining the absolute minimum partial spread intersection parameter $\theta = \min \{s_a, \sigma_{ab}, \sigma_{abc}, \sigma_{1234}\} = 5$.

By Theorem \ref{Thm_Main_Code_Result}, the resulting code $\mathcal{M}_{(\psi_a)}$ is a minimal binary linear code with parameters $[255, 12, 75]$, whose exact weight enumerator polynomial is given by:
\[
\begin{aligned}
W(z) = 1 &+ z^{75} + 5z^{90} + 5z^{105} + 544z^{120} + 675z^{121} + 825z^{122} + 180z^{123} \\
&+ 255z^{128} + 480z^{136} + 525z^{137} + 450z^{138} + 75z^{139}.
\end{aligned}
\] From Table 2, the absolute minimum and maximum Hamming weights evaluate to $w_{\min} = 75$ and $w_{\max} = 140$, respectively. Evaluating their weight ratio gives:
\[
\frac{w_{\min}}{w_{\max}} = \frac{75}{140} > \frac{1}{2}.
\]
Thus, because $\theta = 5 > 2^{m-2} = 4$, this parameter instance directly yields a minimal linear code that \textbf{strictly satisfies the classical Ashikhmin--Barg sufficient condition} ($\frac{w_{\min}}{w_{\max}} > \frac{1}{2}$), serving as an immediate baseline application of Theorem \ref{Thm_Main_Code_Result}. 

\end{example}


\begin{example}\label{ex:AB_violated}
Let $n = 8$ ($m = 4$), yielding code length $N = 2^8 - 1 = 255$ and dimension $K = n + 4 = 12$. We select $17$ pairwise disjoint $4$-dimensional subspaces $\{V_1, V_2, \ldots, V_{17}\}$ in $\mathbb{F}_2^8$. To ensure conditions C1--C3 and the size bounds $2 \le s_a \le 2^{m-1} = 8$ are strictly satisfied, we define the four basis index sets as:
\[
\mathcal{S}_1 = \{1, 2, 3, 4, 10, 11\}, \quad \mathcal{S}_2 = \{3, 4, 5, 6, 10, 11\},
\]
\[
\mathcal{S}_3 = \{1, 3, 7, 8, 10, 12\}, \quad \text{and} \quad \mathcal{S}_4 = \{2, 3, 5, 9, 11, 12\},
\]
giving individual set cardinalities $s_1 = s_2 = s_3 = s_4 = 6 \le 2^{4-1}$.

Evaluating all $15$ non-zero function combinations in $\mathcal{F}$ yields the parameter set:
\[
u \in \{\sigma_{ab}, \sigma_{abc}, \sigma_{1234}\} \subset \{4, 5, 6, 7, 8\},
\]
thereby determining the absolute minimum partial spread intersection parameter $\theta = \min \{s_a, \sigma_{ab}, \sigma_{abc}, \sigma_{1234}\} = 4$ (attained at $\sigma_{12} = 4$).

By Theorem \ref{Thm_Main_Code_Result}, the resulting code $\mathcal{M}_{(\psi_a)}$ is a minimal binary linear code with parameters $[255, 12, 60]$, whose exact weight enumerator polynomial is given by:
\[
\begin{aligned}
W(z) = 1 &+ z^{60} + 5z^{90} + 5z^{105} + 544z^{120} + 750z^{121} + 825z^{122} + 195z^{124} \\
&+ 255z^{128} + 480z^{136} + 525z^{137} + 450z^{138} + 60z^{140}.
\end{aligned}
\]
Here, the absolute minimum and maximum Hamming weights evaluate to $w_{\min} = 60$ and $w_{\max} = 140$, respectively. Evaluating their weight ratio gives:
\[
\frac{w_{\min}}{w_{\max}} = \frac{60}{140}< \frac{1}{2}.
\]
Thus, because $\theta = 4 \le 2^{m-2} = 4$, this parameter instance directly demonstrates a minimal linear code that \textbf{strictly violates the classical Ashikhmin--Barg sufficient condition} ($\frac{w_{\min}}{w_{\max}} \le \frac{1}{2}$), highlighting the geometric sieving capability of Theorem \ref{Thm_Main_Code_Result}.
\end{example}







\section{Applications in Cryptographic Secret Sharing Schemes}

Let $\mathcal{M}$ be an $[N, K, D]$ binary linear code with generator matrix $\mathbf{G} = [\mathbf{g}_0, \mathbf{g}_1, \ldots, \mathbf{g}_{N-1}] \in \mathbb{F}_2^{K \times N}$. In a Massey secret sharing scheme, a trusted dealer distributes a secret $s_0 = \mathbf{u} \cdot \mathbf{g}_0 \in \mathbb{F}_2$ across $N-1$ participants $\mathcal{P} = \{P_1, P_2, \ldots, P_{N-1}\}$ by assigning shares $s_i = \mathbf{u} \cdot \mathbf{g}_i$ for $1 \le i \le N-1$, where $\mathbf{u} \in \mathbb{F}_2^K$ is chosen uniformly at random.

A participant coalition $\mathcal{P}_A = \{P_{i_1}, P_{i_2}, \ldots, P_{i_m}\} \subseteq \mathcal{P}$ is called a \emph{minimal access set} if and only if $\text{rank}([\mathbf{g}_0, \mathbf{g}_{i_1}, \ldots, \mathbf{g}_{i_m}]) = \text{rank}([\mathbf{g}_{i_1}, \ldots, \mathbf{g}_{i_m}])$, with no proper subset satisfying this rank condition. Let $\Gamma$ denote the collection of all minimal access sets. By Massey \cite{Massey1993}, $\mathcal{P}_A \in \Gamma$ if and only if there exists a minimal codeword $\mathbf{c} = (c_0, c_1, \ldots, c_{N-1}) \in \mathcal{M}^\perp$ such that $c_0 = 1$ and $\mathrm{Supp}(\mathbf{c}) = \{0\} \cup \{i_j \mid P_{i_j} \in \mathcal{P}_A\}$.

When $\mathcal{M}$ is a minimal linear code, every non-zero codeword is minimal, guaranteeing that the dual code $\mathcal{M}^\perp$ induces an \emph{ideal} secret sharing scheme where $\Gamma$ is completely characterized by the affine cross-section $c_0 = 1$.



\subsection{Cryptographic Significance of the $n+4$ Access Architecture}
 In this section, we apply the newly constructed family of minimal binary linear codes $\mathcal{M}_{(\psi_a)}$ of dimension $K = n+4$ to design an ideal secret sharing scheme using Massey's foundational framework \cite{Massey1993}.

To formally establish the structural foundation of the access structure induced by our $K = n+4$ code family, we first determine the exact cardinality and support identity of the authorized minimal access sets.


\begin{theorem}\label{Thm:SSS_Construction}
Let $\mathcal{M}_{(\psi_a)}$ be the $[2^n-1, n+4, \theta(2^m-1)]$ minimal binary linear code constructed in Theorem \ref{Thm_Main_Code_Result} over length $N = 2^n - 1$ and dimension $K = n+4$. The access structure $\Gamma$ over $\mathcal{P} = \{P_1, \ldots, P_{2^n-2}\}$ satisfies:
\begin{enumerate}
    \item  $|\Gamma| = 2^{n+3}$.
    \item  For every $\mathcal{P}_A \in \Gamma$, the cardinality $|\mathcal{P}_A|$ is uniquely given by:
    \[
    |\mathcal{P}_A| = w - 1, \quad \text{for some } w \in \Omega \setminus \{0\},
    \]
    where $\Omega$ is the weight spectrum of $\mathcal{M}_{(\psi_a)}$ established in Table 2.
\end{enumerate}
\end{theorem}

\begin{proof}
(i) Define the affine linear subspace $\mathcal{M}_1 = \{\mathbf{c} \in \mathcal{M}_{(\psi_a)} \mid c_0 = 1\}$. Since $\mathcal{M}_{(\psi_a)}$ contains non-zero codewords with non-zero first coordinate, $|\mathcal{M}_1| = 2^{K-1} = 2^{(n+4)-1} = 2^{n+3}$. By Theorem \ref{Thm_Main_Code_Result}, every $\mathbf{c} \in \mathcal{M}_{(\psi_a)} \setminus \{\mathbf{0}\}$ is minimal. Thus, every $\mathbf{c} \in \mathcal{M}_1$ uniquely induces a minimal access set $\mathcal{P}_A = \{P_i \mid c_i = 1, i \neq 0\} \in \Gamma$, establishing $|\Gamma| = 2^{n+3}$.

(ii) For any $\mathbf{c} \in \mathcal{M}_1$, its Hamming weight is $wt(\mathbf{c}) = |\mathrm{Supp}(\mathbf{c})| = 1 + |\mathcal{P}_A|$. Thus, $|\mathcal{P}_A| = wt(\mathbf{c}) - 1$. The values $wt(\mathbf{c})$ are restricted to the non-zero weight support $\Omega \setminus \{0\}$ of $\mathcal{M}_{(\psi_a)}$.
\end{proof}


\begin{thebibliography}{99}
\bibitem[LiuLiao2022]{LiuLiao2022} H.~Liu and Q.~Liao, \textit{A New Constructions of Minimal Binary Linear Codes}, 2022. \url{https://arxiv.org/pdf/2201.02981}
\bibitem[Massey1993]{Massey1993} J.~L.~Massey, \textit{Minimal codewords and secret sharing}, in Proc. 6th Joint Swedish-Russian Int. Workshop on Inf. Theory, M\"olle, Sweden, pp. 276--279, 1993.
\end{thebibliography}
\end{document}
